\documentclass[11pt]{article}
\usepackage[margin=1in]{geometry}
\usepackage{amsmath,amssymb,amsthm}
\usepackage{booktabs}
\usepackage{enumitem}
\usepackage[hidelinks]{hyperref}

\newtheorem{theorem}{Theorem}
\newtheorem{lemma}{Lemma}
\newtheorem{proposition}{Proposition}
\theoremstyle{definition}
\newtheorem{definition}{Definition}
\newtheorem{remark}{Remark}

\newcommand{\C}{\mathbb{C}}
\newcommand{\Q}{\mathbb{Q}}
\newcommand{\Z}{\mathbb{Z}}
\newcommand{\Bhat}{\widehat{\mathcal B}}
\newcommand{\CS}{\operatorname{CS}}
\newcommand{\Vol}{\operatorname{Vol}}
\newcommand{\SL}{\mathrm{SL}}
\newcommand{\PSL}{\mathrm{PSL}}
\newcommand{\Kind}{K_3^{\mathrm{ind}}}

\title{The $7_2$ regulator integral is $4\pi^2/85$}
\author{Rob Sneiderman\\ \texttt{robbysneiderman@gmail.com}}
\date{}

\begin{document}
\maketitle

\begin{abstract}
Problem~3.1 of the Ramanujan Challenge asks for the exact value of a
Godbillon--Vey type knot invariant: an integral of the Chern--Simons
variation form along a distinguished real arc of the $\SL(2,\C)$ character
variety of the prime knot $7_2$.  We give a computer-assisted proof
that
\[
 \int_{\alpha}^{\beta}
 \left(\log x\,\frac{dy}{y}-\log y\,\frac{dx}{x}\right)
 =\frac{4\pi^2}{85}.
\]
The strategy is to realize the whole arc, not just its endpoints, as an exact
algebraic family of $\SL(2,\C)$ representations, and then to recognize the
integral as a difference of two \emph{regulator} values --- Cheeger--Chern--Simons
invariants of the two closed manifolds obtained by Dehn filling at the ends of
the arc.  Concretely: we build an exact continuous real Riley and
enhanced-Ptolemy lift of the full challenge branch; its left endpoint factors
through the $(-1,2)$ Dehn filling of the $7_2$ exterior and its right endpoint
through the $(-1,1)$ filling.  Four exact algebraic tetrahedron shapes at each
endpoint give Neumann's extended Rogers sums $S_\alpha,S_\beta$.  A reciprocal
Galois symmetry of the endpoint number fields, matched to an orientation-preserving
symmetry of the four-tetrahedron triangulation, forces twice each extended Bloch
class into a \emph{totally real} number field; by Borel's rank theorem the
relevant $K$-group is then finite, so each regulator value lies on an explicit
finite lattice.  A $1000$-bit Arb interval evaluation on independently
reconstructed algebraic shapes, with every logarithmic edge
and Dehn-filling equation checked, is much narrower than one lattice spacing and
therefore selects the exact values
\[
 S_\alpha=-\frac{26\pi^2}{15},\qquad
 S_\beta=-\frac{86\pi^2}{51}
 \quad\pmod{\pi^2/2}.
\]
Zickert's Dehn-invariant formula and the differential of the extended Rogers
dilogarithm identify the difference $S_\beta-S_\alpha$ with the required
integral modulo $\pi^2/2$, and a coarse positive bound $0<I<\pi^2/2$ removes the
ambiguity in the lift.  Direct quadrature and a separate signed-coordinate
evaluation reproduce the result to about $60$ and $180$ decimal places,
respectively, but neither numerical agreement is used as proof.
\end{abstract}

\section*{Nature of the computer-assisted proof}

This is a \emph{computer-assisted} proof, and we want to be explicit about what
that means here.  The closing identity $\int_\alpha^\beta\eta=4\pi^2/85$ (here
$\eta=\log M\,d\log L-\log L\,d\log M$ is the Chern--Simons variation form
defined in eq.~\eqref{eq:eta} below) is
verified numerically to about $180$ decimal places; the logical proof does not
rest on that numerical agreement but uses it only as a cross-check.  The proof
proper reduces the integral to a congruence on a finite, explicitly bounded
lattice and then pins down the single lattice point by a certified
interval computation.  We have since independently re-derived the load-bearing pieces of this argument
outside the original certificate's environment.  The whole-arc lift, the single
rational formula $u(M,L)$ of Section~\ref{sec:endpoints} that realizes the entire
branch as one Riley family, was re-derived from scratch in plain \texttt{sympy};
the relator identity of Section~\ref{sec:endpoints}, the enhanced-Ptolemy solution
of Section~\ref{sec:ptolemy}, and the Dehn wedge $\nu=2\,M\wedge L$ of
Section~\ref{sec:variation} were each reproduced, symbolically or to high
precision, from the printed equations alone; and the challenge integral itself was
confirmed to sixty digits by a direct quadrature that touches none of this
machinery.  The two Chern--Simons regulator values $S_\alpha,S_\beta$ were then
reproduced by an independent dilogarithm evaluator to about $150$ digits, fed only
the flattenings and gluing matrix read from SnapPy's public triangulation
interface, so no step rests on the certificate's own Chern--Simons routine.  Those two
topological inputs have since been reconstructed by hand from the knot's
four-tetrahedron triangulation, its explicit face pairings, using no private
extended-Ptolemy solver: the shapes, flattenings, and gluing matrix are recovered
from the gluing equations directly and reproduce the regulator values to over
$160$ digits.  A separate $1000$-bit Arb computation encloses the normalized
regulator difference in a ball of diameter about $5\times10^{-300}$ and selects
$4/85$ as the unique point of the proved regulator lattice.  The only residual
appeal to SnapPy is to read public census combinatorics, the face pairings and
the peripheral gluing rows, which the challenge permits.

Khoi's 2008 paper supplies the general Godbillon--Vey/$A$-polynomial strategy
and evaluates related paths for the figure-eight and $5_2$ knots, but it does
not contain this $7_2$ evaluation.  The formula was presented as an open
conjecture in the 2026 challenge.  This manuscript was submitted to the
organizers on July 12, 2026; organizer review remains pending.

\section{Notation and setup}\label{sec:notation}

This section fixes every symbol used later and gives, for each, a word about
why it appears.  A reader outside low-dimensional topology should be able to
read the rest of the paper after this section.  Two pieces of standard
machinery --- the extended Bloch group and the extended Rogers dilogarithm ---
are recalled where they are first needed, in Sections~\ref{sec:torsion}
and~\ref{sec:variation}.

\paragraph{The curve and the two eigenvalues.}
The knot $7_2$ has a well-defined $\SL(2,\C)$ \emph{character variety}
$X=\operatorname{Hom}(\pi_1(S^3\setminus 7_2),\SL(2,\C))/\!/\SL(2,\C)$, the space of
representations of the knot group up to conjugacy, whose coordinate ring is
generated by trace functions.  Restricting a representation to the boundary
torus and recording the meridian and longitude \emph{eigenvalues} $(M,L)$ (not
their traces) sends the one-dimensional part of $X$ to a plane curve in the
$(M,L)$-plane; the defining polynomial of the Zariski closure of that curve is
the knot's \emph{$A$-polynomial} $A_{7_2}(M,L)$.  The problem statement gives this polynomial explicitly as a
degree-$5$ polynomial in $L$ with coefficients in $\Z[M]$:
\begin{align}
 A_{7_2}(M,L)={}& L^5
 + L^4\bigl(M^{14}-M^{12}+3M^4+4M^2-2\bigr)\notag\\
 &{}+ L^3\bigl(-2M^{18}+5M^{16}+M^{14}-4M^{12}+6M^8+5M^6+2M^4-4M^2+1\bigr)\notag\\
 &{}+ L^2\bigl(M^{22}-4M^{20}+2M^{18}+5M^{16}+6M^{14}-4M^{10}+M^8+5M^6-2M^4\bigr)\notag\\
 &{}+ L\bigl(-2M^{22}+4M^{20}+3M^{18}-M^{10}+M^8\bigr)
 + M^{22}.\label{eq:Apoly}
\end{align}
Here $M$ and $L$ are the two coordinates on the curve, and they have a concrete
meaning for a representation $\rho\colon\pi_1(S^3\setminus 7_2)\to\SL(2,\C)$.
The boundary of the knot exterior is a torus with two distinguished loops, the
\emph{meridian} and the \emph{longitude}.  Their holonomies (the image matrices $\rho(\text{meridian})$
and $\rho(\text{longitude})$ under $\rho$) can be simultaneously
upper-triangularized.  This is not automatic for an arbitrary pair of matrices,
but here the meridian and longitude are the two generators of the fundamental
group $\Z^2$ of the boundary torus, so they commute in $\pi_1$; their holonomies
therefore commute, and a commuting pair of matrices over $\C$ shares an
eigenvector and can be conjugated into upper-triangular form simultaneously.
(The two Riley meridian generators $a,b$ of Section~\ref{sec:endpoints}, by
contrast, do \emph{not} share an eigenvector: that failure is exactly the
irreducibility of $\rho$.)  In this triangular form,
\begin{itemize}[nosep]
 \item $M$ is the \emph{meridian-eigenvalue coordinate}: the function assigning
 to each representation $\rho$ in the family an eigenvalue $M(\rho)$ of the
 meridian holonomy.  At a single fixed $\rho$ this is a definite number, but as
 $\rho$ ranges over the one-dimensional character variety $M$ varies, and it is
 this variation that the integrand records.  It plays the role of the
 $x$-coordinate, $M=x$.
 \item $L$ is the \emph{longitude eigenvalue}: the corresponding eigenvalue of
 the longitude holonomy, in the challenge (= SnapPy standard-longitude)
 normalization.  It is the $y$-coordinate, $L=y$.
\end{itemize}
The vanishing $A_{7_2}(M,L)=0$ is the condition that a pair $(M,L)$ of
eigenvalues actually comes from a representation.  This is the defining property
of the $A$-polynomial~\cite{CCGLS}: $A_{7_2}$ is (the nonabelian factor of) the
defining polynomial of the Zariski closure of the eigenvalue pairs $(M,L)$
realized by the peripheral holonomy of an $\SL(2,\C)$ representation.  We do not
use this abstractly, and in particular do not rely on a point of
$\{A_{7_2}=0\}$ automatically lifting to a representation: the direction we
actually need, that $A_{7_2}(M,L)=0$ forces a genuine representation on the
challenge branch, is established constructively in
Section~\ref{sec:endpoints}, where the Riley relator
identity~\eqref{eq:generic-relator} exhibits an explicit irreducible
representation at every point of the branch.

\paragraph{The Chern--Simons variation form $\eta$.}
The integrand of the challenge is the $1$-form
\begin{equation}\label{eq:eta}
 \eta=\log M\,d\log L-\log L\,d\log M
 =\log x\,\frac{dy}{y}-\log y\,\frac{dx}{x}.
\end{equation}
This is the real $A$-polynomial Chern--Simons variation form studied by
Khoi~\cite{Khoi}; the geometric content of the whole paper is that its integral
along the arc below is a difference of Chern--Simons invariants.  We write $x,y$
and $M,L$ interchangeably as dictated by context: $x,y$ for the curve/integrand,
$M,L$ for the representation-theoretic eigenvalues.

\paragraph{The challenge branch and its endpoints.}
Among the real points of the curve there is one distinguished arc.

\begin{definition}[challenge branch]\label{def:branch}
The \emph{challenge branch} is the real branch $L=y(M)$ of $A_{7_2}(M,L)=0$ for
$\alpha\le M\le\beta$ on which $y$ is positive and decreasing with slope
$y'(x)\le-2$.  For each $x$ in $[\alpha,\beta]$ the quintic in $L$ has five real
roots; the challenge branch is the second-largest of them.  The endpoints are
\[
 \alpha\approx 0.349269,\qquad \beta\approx 0.406813,
\]
where $\alpha$ is the real root of $A_{7_2}(\alpha,\sqrt\alpha)=0$ closest to
$0.349269$ and $\beta$ is the real root of $A_{7_2}(\beta,\beta)=0$ closest to
$0.406813$.  The branch runs from $(M,L)=(\alpha,\sqrt\alpha)$ to $(\beta,\beta)$.
\end{definition}

\noindent To high precision (50 digits),
\[
 \alpha = 0.34926850479999612520833503127884956848473645319259,
\]
\[
 \beta  = 0.40681308133678976238401142887341660342177298944918,
\]
\[
 \tfrac{4\pi^2}{85} = 0.46445197181596981735691722352358358283829173681133.
\]

\paragraph{Endpoint fields, and the roots $t$ and $r$.}
The endpoints are algebraic, and we will need the number fields they generate.
Call these the \emph{endpoint fields}.  They are the number fields generated by
a real root of the two minimal polynomials
\begin{align}
 f_\alpha(t)={}&t^{12}-3t^{11}+4t^{10}-5t^9+6t^8-7t^7+7t^6
 -7t^5+6t^4-5t^3+4t^2-3t+1,\label{eq:fa}\\
 f_\beta(r)={}&r^{16}-7r^{15}+22r^{14}-48r^{13}+87r^{12}-133r^{11}
 +178r^{10}-211r^9+223r^8\notag\\
 &{}-211r^7+178r^6-133r^5+87r^4-48r^3+22r^2-7r+1.\label{eq:fb}
\end{align}
The designated real roots are, plainly,
\begin{equation}\label{eq:tr}
 t=\sqrt\alpha=0.59098942867025644049\ldots,
 \qquad
 r=\beta=0.40681308133678976238\ldots.
\end{equation}
That is: $t$ is exactly the square root of the left endpoint's $M$-value, and
$r$ is exactly the right endpoint's $M$-value.  The reason $t=\sqrt\alpha$ (and
not $\alpha$ itself) is the asymmetric normalization in
Definition~\ref{def:branch}: the left endpoint is defined by $L=M^{1/2}$, so its
longitude eigenvalue is $\sqrt\alpha=t$ and its meridian eigenvalue is
$\alpha=t^2$.  Exact substitution gives the endpoint eigenvalue pairs
\begin{equation}\label{eq:endpts}
 (M_\alpha,L_\alpha)=(t^2,t),\qquad
 (M_\beta,L_\beta)=(r,r),\qquad
 A_{7_2}(M_j,L_j)=0,
\end{equation}
where throughout, the subscript $j\in\{\alpha,\beta\}$ ranges over the two
endpoints and $(M_j,L_j)$ denotes the corresponding meridian/longitude
eigenvalue pair.  We write $f_\alpha$, $f_\beta$ for the polynomials
\eqref{eq:fa}--\eqref{eq:fb} and $E_\alpha=\Q(t)$, $E_\beta=\Q(r)$ for the
endpoint fields they generate.

\paragraph{Riley representations.}
The knot $7_2$ is the two-bridge knot $b(11,5)$ (Schubert normal form).  For a
two-bridge knot, every irreducible $\SL(2,\C)$ representation of the knot group
is conjugate to a two-parameter \emph{Riley representation}: fix one meridian in
upper-triangular form and normalize the conjugate meridian, obtaining generators
\begin{equation}\label{eq:riley}
 a=\begin{pmatrix}s&1\\0&s^{-1}\end{pmatrix},\qquad
 b=\begin{pmatrix}s&0\\-u&s^{-1}\end{pmatrix},
\end{equation}
where $s=M$ is the meridian eigenvalue and $u$ is the \emph{Riley parameter},
the single off-diagonal unknown chosen so that $(a,b)$ satisfies the knot-group
relation.  Fixing the meridian eigenvalue $s=M$, that relation becomes the Riley
polynomial $\phi(M,u)=0$, which has finitely many roots $u$; each root gives a
representation with its own longitude eigenvalue $L$, so along the $A$-polynomial
curve $A_{7_2}(M,L)=0$ exactly one value of $u$ is selected at each point
$(M,L)$.  Thus $u$ is a single-valued rational function $u=u(M,L)$ of the
peripheral eigenvalue pair, given explicitly in \eqref{eq:generic-u} below.
(Warning: the letter $b$ is reused in Section~\ref{sec:ptolemy} for
an unrelated Ptolemy variable.)  The relation and the longitude are built from
the relator word
\begin{equation}\label{eq:word}
 w=aba^{-1}b^{-1}aba^{-1}b^{-1}ab
 \quad(\text{10 letters}),
 \qquad \lambda=w^*wa^{-4},
\end{equation}
where $w^*$ is $w$ read backwards.  The knot-group relation is $wa=bw$, and
$\lambda$ is the canonical longitude of this two-bridge presentation.  The
challenge longitude is the \emph{inverse} of $\lambda$, so if $\lambda_{11}$
denotes the $(1,1)$ entry (an eigenvalue) of the longitude matrix, then the
challenge longitude eigenvalue is
\begin{equation}\label{eq:Llambda}
 L=\lambda_{11}^{-1},\qquad\text{equivalently } L\lambda_{11}=1.
\end{equation}
The identity $L\lambda_{11}=1$ is the ``basis bridge'' between the two-bridge
presentation and the SnapPy standard basis, and we use it again in
Section~\ref{sec:variation}.

\paragraph{Symbols introduced later.}
For orientation, the remaining notation is local to its section and defined
there: the enhanced Ptolemy coordinates ($c_{1100,0}$, $m=M^{-1}$, and the
variables $b,c,e$) and cross-ratios $z_{j,0},\dots,z_{j,3}$ in
Section~\ref{sec:ptolemy}; the shape parameters $z_0,\dots,z_3$, the reciprocal
involutions $\tau$, the fixed-field generators $X=t+t^{-1}$ resp.\ $r+r^{-1}$,
the fixed fields $F_j$, the extended Bloch classes $x_j$, the triangulation
automorphism $\phi$, and the torsion order $w(F)$ in Section~\ref{sec:torsion};
the logarithmic coordinates $w_0,w_1,w_2$ and the regulator sums $S_\alpha,S_\beta$
in Section~\ref{sec:interval}; and the extended Rogers function $R(z)$, the Dehn
invariant $\nu$, and the log-coordinates $X=-\log x$, $Y=-\log y$ in
Section~\ref{sec:variation}.  A handful of letters are deliberately reused
(the matrix $I$ vs.\ the integral $I$; the Riley $b$ vs.\ the Ptolemy $b$; the
polynomial $R(s,u)$ vs.\ the Rogers function $R(z)$; $X$ as a field generator vs.\
$X=-\log x$; the local exponent $\nu_p$ vs.\ the Dehn invariant $\nu$); each
reuse is flagged at its second appearance.

\section{The curve and its two representations}\label{sec:endpoints}

The goal of this section is to promote the two \emph{endpoints} of the challenge
branch to genuine $\SL(2,\C)$ representations, and in fact to lift the
\emph{entire} arc as one algebraic family.  Recognizing the endpoints as
representations of closed Dehn fillings is what later lets us treat the integral
as a difference of Chern--Simons invariants.  Everything here is done by exact
rational-function arithmetic, independently of any triangulation.

\paragraph{One rational formula for the whole arc.}
The key point --- and the step flagged in the honesty note --- is that a single
rational function of $(M,L)$ gives the Riley parameter along the whole branch.
Set $s=M$ and
\begin{equation}\label{eq:generic-u}
 u(M,L)=\frac{(M^2-1)(L-1)}{M^2+L}.
\end{equation}
With this choice, direct rational-function arithmetic in the Riley matrices
\eqref{eq:riley} evaluates the relator $wa-bw$ explicitly.  Its $(1,2)$ entry is
\begin{equation}\label{eq:generic-relator}
 (wa-bw)_{12}=\frac{A_{7_2}(M,L)}{M^6(M^2+L)^5},
 \qquad
 A_{7_2}(M,L)\mid\operatorname{num}(L\lambda_{11}-1)
 \quad\text{in }\Q[M,L],
\end{equation}
and the other three relator entries either vanish identically or are multiples
of this one.  In particular, on the curve $A_{7_2}(M,L)=0$ the relation $wa=bw$
holds, so \eqref{eq:generic-u} really does produce a representation for every
point of the branch.  The same arithmetic yields the commutator trace
\begin{equation}\label{eq:generic-irreducible}
 \operatorname{tr}([a,b])-2=
 \frac{(L-1)(M^2-1)^2(L-M^4)}{M^2(M^2+L)^2},
\end{equation}
where $[a,b]=aba^{-1}b^{-1}$.  Since $\operatorname{tr}([a,b])-2\ne0$ exactly
when the representation is irreducible, \eqref{eq:generic-irreducible} is our
irreducibility detector.

\begin{proposition}\label{prop:arc-lift}
On the challenge branch one has $1/4<M<1/2$, $1/4<L<1$, and $L>M^4$.  Hence every
denominator in \eqref{eq:generic-u}--\eqref{eq:generic-irreducible} is nonzero,
and \eqref{eq:generic-irreducible} never vanishes.  Therefore \eqref{eq:generic-u}
defines a continuous real irreducible representation path over the whole
challenge arc, with positive peripheral eigenvalues different from $1$, i.e.\
purely hyperbolic peripheral holonomy.
\end{proposition}

\begin{proof}
The stated inequalities are exact facts about the branch (they follow from the
root isolation recorded in Section~\ref{sec:variation}).  Given them, the
denominators $M^2+L$, $M^2$, and $M^6(M^2+L)^5$ are all nonzero, so
\eqref{eq:generic-u} and the relator identity \eqref{eq:generic-relator} are
valid; and $(L-1)(M^2-1)^2(L-M^4)\ne0$ because $L\ne1$, $M\ne\pm1$, and
$L>M^4$, so \eqref{eq:generic-irreducible} is nonzero and the representation is
irreducible.  The Riley normalization \eqref{eq:riley} is itself elementary:
triangularize one meridian, then irreducibility lets one normalize the conjugate
meridian as shown.  A reducible knot representation would have longitude
eigenvalue $1$, because its diagonal character factors through the
abelianization, on which the longitude is trivial; here $L\ne1$, consistent with
irreducibility.
\end{proof}

\paragraph{Isolating the two endpoint characters and their fillings.}
Let $R(s,u)$ be the numerator of the $(1,2)$ entry of $wa-bw$.  (This $R(s,u)$
is a polynomial; it is unrelated to Neumann's Rogers function $R(z)$ of
Section~\ref{sec:variation}.)  The diagonal entries of $wa-bw$ vanish
identically and its $(2,1)$ entry has numerator $uR(s,u)$.  Running the exact
Euclidean algorithm in the variable $u$ over each endpoint field isolates a
single value of $u$, via linear greatest common divisors:
\begin{align}
 \gcd_u\bigl(R(t^2,u),t\lambda_{11}(t^2,u)-1\bigr)
 ={}&u+2t^{11}-5t^{10}+5t^9-7t^8+9t^7-9t^6\notag\\
 &{}+9t^5-10t^4+7t^3-6t^2+5t-2,\label{eq:ugcd-a}\\
 \gcd_u\bigl(R(r,u),r\lambda_{11}(r,u)-1\bigr)
 ={}&u+r^{15}-7r^{14}+22r^{13}-48r^{12}+87r^{11}\notag\\
 &{}-133r^{10}+178r^9-211r^8+223r^7-211r^6
 +178r^5\notag\\
 &{}-133r^4+87r^3-48r^2+21r-5.\label{eq:ugcd-b}
\end{align}
Substituting the resulting root for $u$ makes \emph{every} entry of $wa-bw$
zero, and moreover gives the exact matrix identities
\begin{equation}\label{eq:fillings}
 a^{-1}\lambda^{-2}=I\quad(\alpha),
 \qquad
 a^{-1}\lambda^{-1}=I\quad(\beta),
\end{equation}
where $I$ is the $2\times2$ identity matrix.  (This $I$ is unrelated to the
integral $I$ of Section~\ref{sec:variation}.)  These are precisely the relations
of the closed manifolds obtained by Dehn surgery.  (Dehn filling glues a solid
torus to the torus boundary of the exterior so that the peripheral curve
$p\cdot(\text{meridian})+q\cdot(\text{longitude})$ becomes contractible; the
resulting closed manifold's group is the knot group with that curve killed,
which is exactly a relation of the form \eqref{eq:fillings}.)  The left
representation
factors through the SnapPy filling $(-1,2)$, and the right through $(-1,1)$.  One
checks directly that the two gcd roots \eqref{eq:ugcd-a}--\eqref{eq:ugcd-b} are
the specializations of the single generic formula \eqref{eq:generic-u} at
$(t^2,t)$ and $(r,r)$; so the endpoint characters are exactly the ends of the
one-parameter family of Proposition~\ref{prop:arc-lift}.  Finally,
\eqref{eq:Llambda} records the basis bridge $L\lambda_{11}=1$: the variable $L$
is the challenge (and SnapPy standard-longitude) eigenvalue, while $\lambda$ is
the canonical longitude of the two-bridge presentation.

\begin{remark}
The whole-arc lift \eqref{eq:generic-u}, together with the identities
\eqref{eq:generic-relator}--\eqref{eq:generic-irreducible}, was independently
re-derived from scratch in a separate computer algebra system; see the note at
the head of the paper.
\end{remark}

\section{Exact enhanced Ptolemy points}\label{sec:ptolemy}

To compute Chern--Simons invariants we need the representations expressed on an
actual ideal triangulation, as tetrahedron shapes.  We use SnapPy's
four-tetrahedron triangulation of the $7_2$ exterior together with Zickert's
\emph{enhanced Ptolemy variety}~\cite{ZickertPtolemy}: an affine variety whose
points parametrize generically decorated $\SL(2,\C)$
representations on a \emph{fixed} triangulation with prescribed meridian and
longitude eigenvalues.  Here a \emph{decoration} (also called a
boundary-decoration) is a choice of horosphere --- equivalently a coset fixing a
lift of the peripheral fixed point --- at each cusp; this is the extra boundary
data the Ptolemy coordinates record.  Its coordinates are $24$ signed numbers, one per
oriented edge-with-sign; from them one reads off the tetrahedron shapes, and
from the shapes the Cheeger--Chern--Simons regulator.  We recover the same
characters as in Section~\ref{sec:endpoints} on this triangulation.

\paragraph{The gauge-fixed system.}
Normalize one enhanced Ptolemy coordinate to $c_{1100,0}=1$ (a gauge choice),
abbreviate $m=M^{-1}$, and eliminate the coordinates identified by the face
pairings.  The enhanced ideal then reduces to a four-equation system in three
unknowns $b,c,e$:
\begin{equation}\label{eq:ptolemy-system}
 \begin{split}
  -Mb-mc-be&=0,\\
  -Lmb^2-m-e&=0,\\
  Lm^4b+Lmbe-c&=0,\\
  Mc^2-Me^2+c&=0.
 \end{split}
\end{equation}
Here $b,c,e$ are Ptolemy coordinate unknowns and have \emph{nothing} to do with
the Riley generator $b$ of \eqref{eq:riley}.  This system lifts the whole real
arc on one fixed triangulation.  On the curve $A_{7_2}(M,L)=0$ the solution is
\begin{equation}\label{eq:generic-ptolemy}
 b^2=\frac{(1-M^2)(L-M^4)}{LM^2(M^2+L)},\qquad
 e=-\frac{Lb^2+1}{M},\qquad
 c=b(Lb^2+1-M^2).
\end{equation}
The certificate uses an explicit rational choice of $b(M,L)$ and verifies all
four equations of \eqref{eq:ptolemy-system} as rational-function identities
modulo $A_{7_2}$.  As the sharpest such check, the difference between its $b^2$
and the displayed $b^2$ of \eqref{eq:generic-ptolemy} is exactly
\[
 \frac{(M^6+L)\,A_{7_2}(M,L)}
 {L^2(M^2-1)^2M^6(M^2+1)^2(M^2+L)^4},
\]
which vanishes on the curve.

\paragraph{Choosing a continuous branch of shapes.}
The displayed $b^2$ in \eqref{eq:generic-ptolemy} is positive on the challenge
arc, and both endpoint solutions have $b<0$, so we take the continuous negative
branch.  Then
\[
 Lb^2+1-M^2=
 \frac{L(1-M^2)(1+M^2)}{M^2(M^2+L)}>0,
\]
so by \eqref{eq:generic-ptolemy} the coordinates $b,c,e$ stay negative and
nonzero along the arc.  Each of the $24$ signed coordinates is a fixed sign
times a monomial in $M,L$ times one of $1,b,c,e$; hence none vanishes, the four
tetrahedron shapes avoid $0$ and $1$, and their logarithms can be continued from
the left endpoint to the right without a branch change.  This continuity is what
makes the endpoint regulators comparable in Section~\ref{sec:interval}.

\paragraph{Exact endpoint solutions and identification with the Riley characters.}
Over $E_\alpha=\Q(t)$ at $(M,L)=(t^2,t)$ and over $E_\beta=\Q(r)$ at
$(M,L)=(r,r)$, the reduced Gr\"obner basis of \eqref{eq:ptolemy-system} is linear
in $b,c,e$.  The certificate solves it exactly, expands all $24$ signed Ptolemy
coordinates, and substitutes them back into the unreduced enhanced Ptolemy ideal;
every residual is the zero element of the relevant number field.  Writing
$z_{j,0},\dots,z_{j,3}$ for the four primary cross-ratios (tetrahedron shape
parameters) obtained from those coordinates at endpoint $j$, direct
multiplication in the number fields gives every edge product and the filled-cusp
product equal to $1$ exactly, and no shape equals $0$ or $1$.  So the Ptolemy
point is a nondegenerate decorated representation with the stated peripheral
eigenvalues, and its three base coordinates match the specializations of the
rational full-curve coordinates \eqref{eq:generic-ptolemy}.

These endpoint Ptolemy representations are the endpoint Riley representations of
Section~\ref{sec:endpoints}.  Indeed $L_j\ne1$, whereas a reducible knot-group
representation has longitude eigenvalue $1$; so the Ptolemy representation is
irreducible, hence conjugate to the Riley form \eqref{eq:riley}, and the linear
gcds \eqref{eq:ugcd-a}--\eqref{eq:ugcd-b} leave exactly one character over each
peripheral point.  This proves equality of the two \emph{characters}, not merely
of their peripheral eigenvalues.

\begin{remark}
The original certificate obtained this system through SnapPy's enhanced
Ptolemy interface.  The public audit no longer treats that interface as an
oracle: one checker verifies the printed equations, while a second reconstructs
the edge equations, shapes, and flattenings from the explicit four-tetrahedron
census triangulation.  The remaining SnapPy inputs are public face-pairing and
peripheral-curve data.  Both checkers are included in the verification bundle.
\end{remark}

\section{Reciprocal symmetry and a finite regulator lattice}\label{sec:torsion}

We now recall just enough $K$-theory to say what a ``regulator'' is here, and
then exploit a symmetry to trap it on a finite lattice.  Concretely, the
regulator of an endpoint representation is the complex number $S_j$ computed in
Section~\ref{sec:interval} as a sum of extended Rogers dilogarithm values over
the four tetrahedron shapes; it is well-defined only modulo the period recorded
in the table below, and the whole proof consists of pinning down that number.
With that picture fixed, the $K$-theory that follows is structure, not
definition.  Neumann's
\emph{extended Bloch group} $\widehat{\mathcal B}(F)$ of a field $F$ is the
receptacle for a Cheeger--Chern--Simons class of a representation together with a
choice of flattening (a lift of tetrahedron shape logarithms); it is naturally
isomorphic to the indecomposable part $\Kind(F)$ of algebraic $K$-theory, and it
carries the extended Rogers regulator (defined in Section~\ref{sec:variation}).
There is an $\SL$ version $\widehat{\mathcal B}_{\SL}(F)$ and a $\PSL$ version,
differing by the period conventions recorded below.  The value $\int\eta$ we
want will be a difference of regulator values, so the whole difficulty is to
know that difference lies in a small, explicitly known group.

\paragraph{Reciprocal involutions and totally real fixed fields.}
Both endpoint polynomials \eqref{eq:fa}--\eqref{eq:fb} are reciprocal, so
$t\mapsto t^{-1}$ and $r\mapsto r^{-1}$ are field automorphisms.  Let
$\tau_\alpha(t)=t^{-1}$ on $E_\alpha$ and $\tau_\beta(r)=r^{-1}$ on $E_\beta$;
write $\tau$ for either one.  Its fixed field is generated by $X=t+t^{-1}$
(resp.\ $X=r+r^{-1}$), with defining polynomials
\begin{align}
 g_\alpha(X)&=X^6-3X^5-2X^4+10X^3-X^2-7X+1,\label{eq:ga}\\
 g_\beta(X)&=X^8-7X^7+14X^6+X^5-25X^4+9X^3+12X^2-3X-1.\label{eq:gb}
\end{align}
(This field generator $X$ is unrelated to $X=-\log x$ in
Section~\ref{sec:variation}.)  Exact real-root isolation gives signatures
$(6,0)$ and $(8,0)$ and discriminants $5^3\cdot3881$ and $17^7$; in particular
both fixed fields, which we call $F_\alpha$ and $F_\beta$, are \emph{totally
real}.

\paragraph{Matching the field symmetry to a symmetry of the triangulation.}
The reciprocal action on the exact tetrahedron shapes is the same at both
endpoints.  Writing $z_0,z_1,z_2,z_3$ for the four generic shape parameters
(recall a shape $z$ has partners $1-z^{-1}$, etc.),
\begin{equation}\label{eq:shape-action}
 \tau(z_0)=1-z_2^{-1},\qquad
 \tau(z_1)=z_1,\qquad
 \tau(z_2)=(1-z_0)^{-1},\qquad
 \tau(z_3)=z_3.
\end{equation}
The four-tetrahedron triangulation has a nonidentity orientation-preserving
automorphism $\phi$,
\begin{equation}\label{eq:tet-auto}
 0\mapsto(2;(2,1,3,0)),\quad
 1\mapsto(1;(3,2,1,0)),\quad
 2\mapsto(0;(3,1,0,2)),\quad
 3\mapsto(3;(1,0,3,2)),
\end{equation}
where the tuple after each semicolon is the vertex permutation.  All four
permutations are even (so $\phi$ preserves orientation), and the induced cusp map
is $-I$, so $\phi$ extends over both fillings.  Its shape transformations are
exactly \eqref{eq:shape-action}.  We write $\phi_*$ for the induced map on
representations.

\paragraph{Twice each Bloch class descends to the totally real field.}
The exact Riley lift is an $\SL_2(E_j)$ representation of the closed filling, so
it defines an extended Bloch class
\[
 x_j\in\widehat{\mathcal B}_{\SL}(E_j)\simeq\Kind(E_j),
\]
with image $\bar x_j$ in the $\PSL$ extended Bloch group.  The two groups and
their periods must be kept separate:
\begin{center}
\begin{tabular}{@{}lll@{}}
\toprule
 group / class & period & reference\\
\midrule
 $\SL$ regulator & mod $4\pi^2$ & \cite[Section~2.2.1]{ZickertK3}\\
 $\PSL$ Cheeger--Chern--Simons & mod $\pi^2$ & \cite[Prop.~2.5, Thm.~2.6]{Neumann}\\
 certified evaluator & mod $\pi^2/2$ & \cite[Thm.~3.4]{Yoon}\\
\bottomrule
\end{tabular}
\end{center}
All maps here are homomorphisms, so an order bound for $x_j$ also bounds the
evaluator's value on $\bar x_j$.  Now the shape identities \eqref{eq:shape-action}
say that $\tau_j\rho_j$ is $\PSL_2$-conjugate to $\rho_j\circ\phi_*$, where
$\rho_j$ is the endpoint representation and $\phi$ is the automorphism
\eqref{eq:tet-auto}.  Nondegenerate shape data determine the decorated $\PSL$
representation; $\phi$ preserves orientation and its cusp map is $-I$, so it
extends over each filling and preserves the closed fundamental class.  The two
resulting $\SL$ lifts can therefore differ only by a $\Z/2$ cocycle, and
Zickert's two-torsion result for such a change of lift gives
\cite[Theorem~6.15]{ZickertK3}
\[
 2\tau_j(x_j)=2x_j.
\]
Galois descent \cite[Corollaries~7.14--7.15]{ZickertK3} then places $2x_j$ in
$\Kind(F_j)$, where $F_j$ is the totally real fixed field of $\tau_j$.

\paragraph{Borel finiteness and the explicit lattice.}
Borel's rank theorem states $\operatorname{rank}\Kind(F)=r_2(F)$, the number of
complex (conjugate-pair) places of $F$~\cite{Borel}.  Since $F_j$ is totally
real, $r_2(F_j)=0$, so $\Kind(F_j)$ is \emph{finite}.  Its torsion subgroup is
cyclic of order
\begin{equation}\label{eq:wF}
 w(F)=2\prod_p p^{\nu_p},\qquad
 \nu_p=\max\{\nu:\zeta_{p^\nu}+\zeta_{p^\nu}^{-1}\in F\}
\end{equation}
\cite[Section~1]{ZickertK3}, where $\zeta_{p^\nu}$ is a primitive $p^\nu$-th root
of unity and $\zeta+\zeta^{-1}$ is the real generator of its cyclotomic field.
(This local exponent $\nu_p$ is unrelated to the Dehn invariant $\nu$ of
Section~\ref{sec:variation}.)  We need only degree-based upper bounds: if $F$ is
totally real of degree $d=[F:\Q]$ and contains $\zeta_{p^\nu}+\zeta_{p^\nu}^{-1}$,
then $\varphi(p^\nu)/2$ divides $d$ (for $p^\nu>2$), with $\varphi$ Euler's
totient.  Applying this to $F_\alpha$ (degree $6$): the condition
$\varphi(p^\nu)\mid12$ is met only by the prime powers $p^\nu=8$
($\varphi=4$), $9$ ($\varphi=6$), $5$ ($\varphi=4$), $7$ ($\varphi=6$),
$13$ ($\varphi=12$), giving factors $2^3,3^2,5,7,13$; for $F_\beta$ (degree
$8$) the condition $\varphi(p^\nu)\mid16$ is met only by $32$ ($\varphi=16$),
$3$ ($\varphi=2$), $5$ ($\varphi=4$), $17$ ($\varphi=16$), giving
$2^5,3,5,17$.  With the leading factor $2$ of \eqref{eq:wF} this yields
\begin{equation}\label{eq:w-bounds}
 w(F_\alpha)\mid 2\cdot2^3\cdot3^2\cdot5\cdot7\cdot13=65520,
 \qquad
 w(F_\beta)\mid 2\cdot2^5\cdot3\cdot5\cdot17=16320.
\end{equation}
Thus $x_\alpha$ and $x_\beta$ have orders dividing $131040$ and $32640$.  After
projection to $\C/(\pi^2/2)\Z$, the regulator difference has order dividing
\[
 N=\operatorname{lcm}(131040,32640)=8910720,
\]
so, after division by $\pi^2$, it lies on the finite lattice
\begin{equation}\label{eq:lattice}
 \frac{1}{2N}\Z\bigg/\frac12\Z
 =\frac{1}{17821440}\Z\bigg/\frac12\Z.
\end{equation}

\section{Certified selection of the regulator value}\label{sec:interval}

The regulator difference is now known to sit on the lattice \eqref{eq:lattice},
whose spacing is about $5.61\times10^{-8}$.  It remains to identify the exact
lattice point.  We do this by a certified $1000$-bit Arb interval evaluation of the
extended Rogers sum --- the sum of Neumann's extended Rogers function $R(z)$
(defined in Section~\ref{sec:variation}) over the four tetrahedra --- at each
endpoint, narrow enough to pin one point.

\paragraph{Flattenings with no hidden $2\pi i$.}
For each endpoint, the independently reconstructed algebraic shapes determine
logarithmic coordinates $w_0,w_1,w_2$ on every tetrahedron (branch-chosen logs
of the shape parameters).  The audit multiplies these logarithms by the full
$5\times12$ PGL (projective general linear) gluing matrix, whose rows encode the
four edge relations plus the
Dehn-filling relation labelled \texttt{filling\_0\_0}.  All four edge rows and the
filling row evaluate to zero as complex intervals; the largest row diameter is
$2.24\times10^{-207}$ at $\alpha$ and $6.28\times10^{-206}$ at $\beta$.  Because
all peripheral eigenvalues are positive, the filling relations carry no hidden
$2\pi i$ term:
\[
 -\log M_\alpha+2\log L_\alpha=0,
 \qquad
 -\log M_\beta+\log L_\beta=0.
\]
Thus the closed strong flattenings (flattenings whose log-parameters sum to the
value Neumann requires for the closed manifold, a strengthening of the plain
flattening of Section~\ref{sec:torsion}) are the endpoint values of the
continuous exterior Ptolemy regulator of Section~\ref{sec:ptolemy}.

\paragraph{Evaluating the extended Rogers sum.}
The extended Rogers sum is evaluated with $1000$-bit real balls, together
with Neumann's $6$-torsion correction for a triangulation without a vertex
ordering \cite[Section~13 and Proposition~13.1]{Neumann}.  The lifted flattenings
satisfy Yoon's edge and Dehn-filling conditions; his Theorem~3.1 and the parity
argument in the proof of Theorem~3.4 identify the sum with the filled complex
volume modulo $\pi^2/2$~\cite{Yoon}.  The correction coefficient is exactly $-6$
at both endpoints, contributing $-6(\pi^2/6)=-\pi^2$, which is $0$ modulo
$\pi^2/2$ and cancels in the difference.  In the evaluator's $i(\Vol+i\CS)$
convention, where $\Vol$ and $\CS$ denote the hyperbolic volume and
Chern--Simons invariant of the filled manifold, the sums are
\begin{align}
 \frac{S_\alpha}{\pi^2}
  &=-1.733333333333333333333333\ldots,
  \label{eq:Sa}\\
 \frac{S_\beta}{\pi^2}
  &=-1.686274509803921568627451\ldots.
  \label{eq:Sb}
\end{align}

\paragraph{The lattice selects a unique value.}
The Arb enclosure of $(S_\beta-S_\alpha)/\pi^2$ contains $4/85$ and has
diameter less than $5\times10^{-300}$, far below the lattice spacing
$5.61\times10^{-8}$ of \eqref{eq:lattice}.  Both neighboring lattice points
$4/85\pm1/(2N)$ are rigorously excluded.  Hence the enclosure contains exactly
one admissible regulator value, so
\begin{equation}\label{eq:reg-difference}
 S_\beta-S_\alpha=\frac{4\pi^2}{85}
 \quad\pmod{\pi^2/2}.
\end{equation}
The same lattice argument selects the individual representatives
\[
 S_\alpha=-\frac{26\pi^2}{15},\qquad
 S_\beta=-\frac{86\pi^2}{51}
 \quad\pmod{\pi^2/2}
\]
(consistent with \eqref{eq:Sa}--\eqref{eq:Sb}, since $-26/15=-1.7333\ldots$ and
$-86/51=-1.686274\ldots$).

\paragraph{An independent numerical cross-check.}
The original certificate also runs a separate numerical path that does not feed the
lattice proof.  It expands the actual signed Ptolemy coordinates, uses SnapPy's
public cross-ratio and flattening code, and multiplies the logarithmic gluing
matrix explicitly.  This path uses the $\Vol+i\CS$ convention (rather than
$i(\Vol+i\CS)$) and returns
\[
 \frac{\CS_\alpha}{\pi^2}=\frac1{15},\qquad
 \frac{\CS_\beta}{\pi^2}=\frac1{51}
 \quad\pmod{1/6},
\]
so $\CS_\alpha-\CS_\beta=4\pi^2/85$ to about $180$ decimal places (with $\Vol$
and $\CS$ as above, the hyperbolic volume and Chern--Simons invariant).

\section{From the regulator to the challenge integral}\label{sec:variation}

Two things remain: to check that the regulator difference is genuinely the
integral of $\eta$ (not some other $1$-form), and to remove the modulo $\pi^2/2$
ambiguity.  The first is a differential identity built on the Dehn invariant and
the extended Rogers dilogarithm; the second is a coarse but rigorous size bound.

\paragraph{The extended Rogers dilogarithm and the Dehn invariant.}
Let $S(M,L)$ be the continuous extended Rogers sum on the Ptolemy lift
\eqref{eq:generic-ptolemy}, with logarithms continued from the left endpoint; its
endpoint values are the $S_\alpha,S_\beta$ of Section~\ref{sec:interval}.  The
unordered-triangulation correction is constant, so its differential is zero.
Neumann's \emph{extended Rogers function} $R(z)$ (on the extended Bloch group,
unrelated to the polynomial $R(s,u)$ of Section~\ref{sec:endpoints}) has
differential
\begin{equation}\label{eq:dR}
 2\,dR(z)=\log z\,d\log(1-z)-\log(1-z)\,d\log z.
\end{equation}
The remaining ingredient is the Dehn invariant, which measures how the shapes
vary along the peripheral torus.  For $n=2$, Zickert's Dehn-invariant formula
gives $-2\,M\wedge L$ in a peripheral basis of intersection number $+1$
\cite[Theorem~1.10]{ZickertPtolemy}, where $M\wedge L$ is the wedge of the
meridian and longitude eigenvalues in the target of the Dehn invariant.  Our
$M,L$ are the challenge and SnapPy standard-basis eigenvalues by the exact
identity $L\lambda_{11}=1$ from Section~\ref{sec:endpoints}, and the standard
knot meridian--longitude basis has intersection number $-1$
\cite[Section~3.1]{ZickertPtolemy}.  In that basis the Dehn invariant is
therefore
\begin{equation}\label{eq:dehn}
 \nu=2\,M\wedge L.
\end{equation}
As a direct sign check, the independent script
\texttt{dehn\_wedge\_reconstruction.py} forms $z_i\wedge(1-z_i)$ from the
exponent vectors of all $24$ signed-coordinate monomials; every wedge involving
the Ptolemy variables $b,c,e$ cancels, leaving the single surviving coefficient
$+2$ on $M\wedge L$.

\begin{remark}
The sign in \eqref{eq:dehn} is no longer accepted only from the original
certificate.  The independent reconstruction enumerates the orientation signs,
cancels every wedge involving the internal Ptolemy variables, and leaves the
single coefficient $+2$ on $M\wedge L$.
\end{remark}

\paragraph{The differential is exactly $\eta$.}
Applying \eqref{eq:dR} to the formal shape sum and adding the Dehn term
\eqref{eq:dehn} gives
\begin{equation}\label{eq:dS}
 dS=\log M\,d\log L-\log L\,d\log M=\eta.
\end{equation}
Since the continuous Ptolemy lift connects the exact endpoint flattenings,
integrating \eqref{eq:dS} and using \eqref{eq:reg-difference} yields the integral
modulo the period,
\begin{equation}\label{eq:congruence}
 I:=\int_\alpha^\beta\eta=\frac{4\pi^2}{85}
 \quad\pmod{\pi^2/2}.
\end{equation}
(This integral $I$ is unrelated to the identity matrix $I$ of
Section~\ref{sec:endpoints}.)

\paragraph{Removing the period ambiguity.}
The stated branch fixes the real lift with no numerical quadrature.  Put
$X=-\log x$ and $Y=-\log y$, both positive since $x,y\in(0,1)$.  (This $X$ is
unrelated to the fixed-field generator $X$ of Section~\ref{sec:torsion}.)  Since
$dx>0$ and $dy<0$ along the branch,
\[
 \eta=X\,dY-Y\,dX>0,
\]
so $I>0$.  Exact root isolation gives
\[
 \frac14<\alpha<\frac38<\beta<\frac12,
 \qquad
 \frac14<y(\beta)<\frac12<y(\alpha)<1,
\]
with endpoint values $y(\alpha)=\sqrt\alpha\approx0.5909894287$ and
$y(\beta)=\beta\approx0.4068130813$.  It follows that
\begin{align*}
 I
 &<\log4\bigl((Y_\beta-Y_\alpha)+(X_\alpha-X_\beta)\bigr)\\
 &=\log4\,\log\!\left(
 \frac{y(\alpha)}{y(\beta)}\frac{\beta}{\alpha}
 \right)
 <\log4\,\log8<\frac92<\frac{\pi^2}{2}.
\end{align*}
Both $I$ and $4\pi^2/85$ lie strictly between $0$ and $\pi^2/2$, so the
congruence \eqref{eq:congruence} has exactly one admissible real lift.  This
proves the challenge identity:
\[
 \boxed{\displaystyle
 \int_{\alpha}^{\beta}
 \left(\log x\,\frac{dy}{y}-\log y\,\frac{dx}{x}\right)
 =\frac{4\pi^2}{85}.}
\]

\section{Reproducibility}\label{sec:repro}

The public release uses the independent-verification route rather than requiring
SnapPy's private \texttt{dev.extended\_ptolemy} interface.  It requires Python
with \texttt{sympy} and \texttt{mpmath}, plus SageMath~10.9 with SnapPy~3.3.2
installed in Sage's Python.  From the repository root, the complete gate is
\begin{verbatim}
verification/run_all.sh
\end{verbatim}
The runner first checks that Sage and SnapPy share one interpreter, then executes
the exact Riley identities, independent endpoint and uniqueness checks,
Dehn-wedge reconstruction, direct quadrature, by-hand triangulation and
flattening reconstruction, and the final Arb lattice-selection certificate.
The load-bearing final command is
\begin{verbatim}
$SAGE -python verification/exactness_closure.py
\end{verbatim}
and prints its explicit closed verdict only after the field, torsion-bound,
interval, and neighboring-lattice-point assertions pass.  The exact status of
every layer is recorded in the verification ledger shipped with the source.
Borel's finiteness theorem and
the Zickert--Neumann regulator normalizations are cited inputs, not reproved.

\section*{Acknowledgments and tooling disclosure}

The Ramanujan Machine group posed the problem.  The project was developed by
Rob Sneiderman with extensive assistance from OpenAI Codex and other
frontier language models for symbolic exploration, implementation, adversarial
review, and exposition.  Model output is not treated as mathematical evidence:
the claim rests on the displayed argument, cited published theorems, and the
reproducible exact and certified computations described above.

\begin{thebibliography}{9}

\bibitem{CCGLS}
D.~Cooper, M.~Culler, H.~Gillet, D.~D.~Long, P.~B.~Shalen,
\emph{Plane curves associated to character varieties of $3$-manifolds},
Invent. Math. \textbf{118} (1994), no.~1, 47--84.

\bibitem{Borel}
A.~Borel,
\emph{Cohomologie de $\mathrm{SL}_n$ et valeurs de fonctions z\^eta aux points entiers},
Ann. Scuola Norm. Sup. Pisa Cl. Sci. (4) \textbf{4} (1977), no.~4, 613--636.

\bibitem{Khoi}
V.~T.~Khoi,
\emph{On the integral of $\log x\,dy/y-\log y\,dx/x$ over the A-polynomial curves},
Acta Math. Vietnam. \textbf{33} (2008), no.~3, 519--528;
arXiv:0811.2725.

\bibitem{Neumann}
W.~D.~Neumann,
\emph{Extended Bloch group and the Cheeger--Chern--Simons class},
Geom. Topol. \textbf{8} (2004), 413--474;
arXiv:math/0307092.

\bibitem{Yoon}
S.~Yoon,
\emph{The volume and Chern--Simons invariant of a Dehn-filled manifold},
Topology Appl. \textbf{256} (2019), 105--121;
arXiv:1801.08288.

\bibitem{ZickertPtolemy}
C.~K.~Zickert,
\emph{Ptolemy coordinates, Dehn invariant and the A-polynomial},
Math. Z. \textbf{283} (2016), 515--537;
arXiv:1405.0025.

\bibitem{ZickertK3}
C.~K.~Zickert,
\emph{The extended Bloch group and algebraic $K$-theory},
J. Reine Angew. Math. \textbf{704} (2015), 21--54;
arXiv:0910.4005.

\end{thebibliography}

\end{document}
